% Unit 114 English translation of chapter8.tex lines 1353--1482.
% Source: Methods of Algebra, Volume 2, commit
% 9a5803ff77dd3257484cb177f851a73770a59dd3 (CC BY 4.0).
% stable-unit-id: o014.aljabr2.chapter8.bar-construction-b
% Coverage: Section 8.7 from the meaning of a bar resolution through the group-homology example.
% Editorial corrections: the augmented bar complex for a general module is
% C^{aug}(Bar(M)), not C^{aug}(Bar(R)); a missing tensor sign in the final
% group-homology differential is restored from the Indonesian correction witness.

% segment-id: o014.li2.chapter8.bar-construction-b.p001
The augmented simplicial object in
\eqref{eqn:bar-simplicial-resolution} given by the bar construction can be
viewed heuristically as a replacement for $M$. To clarify this idea, consider
the following question. Let $\Gamma$ be a group, and equip the one-point set
$\{\mathrm{pt}\}$ with the trivial right $\Gamma$-action. How should we
understand
% segment-id: o014.li2.chapter8.bar-construction-b.q001
\[ \text{the ``quotient space''}\; \{\mathrm{pt}\} / \Gamma ? \]
With the naive definition of a quotient set or topological space, the answer
is of course uninteresting; in general, taking the quotient by a nonfree
action erases information from the original space. Various needs arising in
practice, however, prompt us to seek a more refined definition of a quotient
by a nonfree action.

% segment-id: o014.li2.chapter8.bar-construction-b.p002
To this end, write $\cated{Set}\Gamma$ for the category of all small right
$\Gamma$-sets, and consider the free--forgetful adjunction
$\cate{Set}\leftrightarrows\cated{Set}\Gamma$. From the one-point right
$\Gamma$-set $\{\mathrm{pt}\}$, the bar construction produces a simplicial
object in $\cated{Set}\Gamma$; this is precisely the simplicial set
$\mathrm{E}\Gamma$ introduced in Example~\ref{eg:classifying-space}. The
detailed verification is left as an exercise in this chapter. Once this fact
is understood and the bar construction is accepted as a replacement for
$\{\mathrm{pt}\}$ in this setting, it is reasonable to define
$\{\mathrm{pt}\}/\Gamma$ as
$(\mathrm{E}\Gamma)/\Gamma\simeq\mathrm{B}\Gamma$. This is a quotient by a
free action and gives the classifying space of $\Gamma$; it is more natural,
and usually more useful, than the naive quotient set.

% segment-id: o014.li2.chapter8.bar-construction-b.p003
In what sense can the augmented simplicial objects in
Examples~\ref{eg:bar-adjunction} and~\ref{eg:bar-comonad} be said to provide a
kind of replacement, or ``resolution'', of the original object? The idea
comes from topology, and its precise meaning involves the contractibility of
Definition~\ref{def:contractible-aug}. One related result is the following.

% segment-id: o014.li2.chapter8.bar-construction-b.pr001
\begin{proposition}\label{prop:bar-contractible}
	In the setting of Example~\ref{eg:bar-adjunction}, suppose we are given an
	adjoint pair
	% segment-id: o014.li2.chapter8.bar-construction-b.q002
	\[\begin{tikzcd}
		F : \mathcal{C} \arrow[shift left, r] & \mathcal{D}: G \arrow[shift left, l] ,
	\end{tikzcd} \quad \begin{array}{ll}
		\eta: \identity_{\mathcal{C}} \to GF & \text{unit} \\
		\epsilon: FG \to \identity_{\mathcal{D}} & \text{counit}
	\end{array}\]
	and form the augmented simplicial object $\mathrm{Bar}(F,G)$ in
	$\End(\mathcal{D})$. Then $G\mathrm{Bar}(F,G)$ is left contractible.
\end{proposition}
\begin{proof}
	Use the notation of Example~\ref{eg:bar-adjunction}, such as $T=GF$ and
	$\mu:T^2\to T$. Recall that
	$G\mathrm{Bar}(F,G)_n=T^{n+1}G$. For every $n\geq-1$, define
	% segment-id: o014.li2.chapter8.bar-construction-b.q003
	\[ k_n := \eta T^{n+1}G: T^{n+1}G \to T^{n+2}G . \]
	Recall also that the face morphism $d_i:T^{n+1}G\to T^nG$ of
	$G\mathrm{Bar}(F,G)$ is $T^i\mu T^{n-i-1}G$ if $0\leq i<n$, and
	$T^nG\epsilon$ if $i=n$; its augmentation morphism is
	$G\epsilon:TG\to G$ (note that this notation differs from that in
	Definition~\ref{def:contractible-aug}). We now verify the conditions for
	left contractibility one by one.

	First, the triangle identity for the adjunction gives
	$(G\epsilon)k_{-1}=(G\epsilon)(\eta G)=\identity_G$.

	Second, $(T,\mu,\eta)$ is a monad, so
	$\mu(\eta T)=\identity_T$. If $n\geq0$, composing both sides on the right
	with $T^nG$ gives
	$d_0k_n=\mu T^nG\cdot\eta T^{n+1}G=\identity_{T^{n+1}G}$.

	Finally, by the naturality of $\eta$, the following diagram commutes when
	$1\leq i\leq n+1$:
	% segment-id: o014.li2.chapter8.bar-construction-b.q004
	\[\begin{tikzcd}[column sep=large]
		T^{n+1} G \arrow[r, "{k_n = \eta T^{n+1}G}" inner sep=0.6em] \arrow[d, "{d_{i-1} = \varphi}"'] & T^{n+2}G \arrow[d, "{d_i = T\varphi}"] \\
		T^n G \arrow[r, "{k_{n-1} = \eta T^n G}"' inner sep=0.6em] & T^{n+1} G
	\end{tikzcd} \quad
		\varphi := \begin{cases}
			T^{i-1} \mu T^{n-i} G, & 1 \leq i \leq n \\
			T^n G\epsilon, & i = n+1.
		\end{cases}\]
	This gives $d_1k_0=k_{-1}(G\epsilon)$ when $i=1$ and $n=0$, and
	$d_ik_n=k_{n-1}d_{i-1}$ when $n\geq1$. The result follows.
\end{proof}

% segment-id: o014.li2.chapter8.bar-construction-b.p004
Although the contractibility of $G\mathrm{Bar}(F,G)$ does not necessarily
imply that $\mathrm{Bar}(F,G)$ is contractible,
Proposition~\ref{prop:bar-contractible} is already sufficient for the
applications below. The exercises in this chapter will introduce further
results concerning contractibility.

% segment-id: o014.li2.chapter8.bar-construction-b.p005
For additive categories, the meaning of ``resolution'' can also be carried to
the level of chain complexes and thus understood through linear algebra.
First, by Definition~\ref{def:augmented-C-cplx}, an augmented simplicial
object $X$ in an additive category gives an augmentation
$\mathrm{C}^{\mathrm{aug}}X$ of the chain complex $\mathrm{C}X$. If the
identity morphism of a chain complex is null-homotopic, the chain complex
itself is called \emph{null-homotopic}.

% segment-id: o014.li2.chapter8.bar-construction-b.th001
\begin{theorem}\label{prop:bar-null-homotopic}
	Let $\mathcal{A}$ be an additive category, suppose an adjoint pair is given,
	% segment-id: o014.li2.chapter8.bar-construction-b.q005
	\[\begin{tikzcd}
		F : \mathcal{A} \arrow[shift left, r] & \mathcal{B}: G, \arrow[shift left, l]
	\end{tikzcd}\]
	and form the augmented simplicial object $\mathrm{Bar}(F,G)$ in
	$\End(\mathcal{B})$ (respectively, $G\mathrm{Bar}(F,G)$ in
	$\mathcal{A}^{\mathcal{B}}$).
	\begin{itemize}
		\item Notice that $\mathcal{A}^{\mathcal{B}}$ is an additive category
		(see the first part of \S\ref{sec:Abel-cat-def}). The chain complex
		$\mathrm{C}^{\mathrm{aug}}\bigl(G\mathrm{Bar}(F,G)\bigr)$ from
		Definition~\ref{def:augmented-C-cplx} is null-homotopic.
		\item For every $Y\in\Obj(\mathcal{B})$, evaluate
		$\mathrm{Bar}(F,G)$ (respectively, $G\mathrm{Bar}(F,G)$) at $Y$
		degreewise. This produces an augmented simplicial object
		$\mathrm{Bar}(Y)$ in $\mathcal{B}$ (respectively,
		$G\mathrm{Bar}(Y)$ in $\mathcal{A}$), and the chain complex
		$\mathrm{C}^{\mathrm{aug}}\bigl(G\mathrm{Bar}(Y)\bigr)$ in
		$\mathcal{A}$ is null-homotopic.
	\end{itemize}
\end{theorem}
\begin{proof}
	Proposition~\ref{prop:bar-contractible} shows that
	$G\mathrm{Bar}(F,G)$ is left contractible. Since contractibility is an
	absolute condition, $G\mathrm{Bar}(Y)$ obtained through the evaluation
	functor $\mathrm{ev}_Y:\mathcal{A}^{\mathcal{B}}\to\mathcal{A}$ is
	naturally left contractible as well. Apply
	Proposition~\ref{prop:contractible-null-homotopic} to both objects.
\end{proof}

% segment-id: o014.li2.chapter8.bar-construction-b.p006
Now suppose that $\mathcal{A}$ and $\mathcal{B}$ are both additive
categories. For an adjoint pair as in
Theorem~\ref{prop:bar-null-homotopic} and $Y\in\Obj(\mathcal{B})$, since
$\mathrm{Bar}(Y)_{-1}=Y$, flattening
$\mathrm{C}^{\mathrm{aug}}(\mathrm{Bar}(Y))$ gives a morphism in
$\cate{Ch}_{\geq0}(\mathcal{B})$
% segment-id: o014.li2.chapter8.bar-construction-b.q006
\begin{equation}\label{eqn:bar-resolution}
	\epsilon_Y: \mathrm{C}(\mathrm{Bar}(Y)) \to Y;
\end{equation}
it is canonical in $Y$, and its degree-zero part is precisely the augmentation
morphism $\epsilon_Y:FG(Y)\to Y$ of $\mathrm{Bar}(Y)$, so the notation is
appropriate.

% segment-id: o014.li2.chapter8.bar-construction-b.co001
\begin{corollary}\label{prop:bar-exactness-G}
	Suppose an adjoint pair of abelian categories is given\footnote{Here $F$
	and $G$ are necessarily additive functors; see
	Corollary~\ref{prop:automatic-additivity}.}
	% segment-id: o014.li2.chapter8.bar-construction-b.q007
	$\begin{tikzcd}
		F : \mathcal{A} \arrow[shift left, r] & \mathcal{B}: G \arrow[shift left, l]
	\end{tikzcd}$
	and let $Y\in\Obj(\mathcal{B})$. The morphism obtained by applying $G$
	degreewise to \eqref{eqn:bar-resolution}, namely
	% segment-id: o014.li2.chapter8.bar-construction-b.q008
	\[ G\epsilon_Y: \mathrm{C}(G\mathrm{Bar}(Y)) \to GY ,\]
	is a quasi-isomorphism in $\cate{Ch}(\mathcal{A})$.
\end{corollary}
\begin{proof}
	Theorem~\ref{prop:bar-null-homotopic} implies that
	$\mathrm{C}^{\mathrm{aug}}(G\mathrm{Bar}(Y))$ is acyclic, and the result
	follows immediately.
\end{proof}

% segment-id: o014.li2.chapter8.bar-construction-b.eg001
\begin{example}[The bar resolution of a module and Hochschild homology/cohomology]\label{eg:HH-comonad}
	\index{bar resolution}
	\index[sym1]{BR}
	\index[sym1]{B'R}
	\index[sym1]{HH}
	Let $\Bbbk$ be a commutative ring and $R$ a $\Bbbk$-algebra. Write
	$\otimes:=\otimes_{\Bbbk}$. By Example~\ref{eg:comonad-change-ring}, the
	comonad determined by the adjoint pair
	% segment-id: o014.li2.chapter8.bar-construction-b.q009
	\[\begin{tikzcd}
		R \otimes (\cdot): \Bbbk\dcate{Mod} \arrow[shift left, r] & R\dcate{Mod}: \text{forgetful} \arrow[shift left, l]
	\end{tikzcd}\]
	is the endofunctor $L:M\mapsto R\otimes M$ on $R\dcate{Mod}$. There is a
	corresponding version for right $R$-modules, which we shall not repeat.

	Now apply the bar construction. For every left $R$-module $M$, we have a
	canonical morphism in $\cate{Ch}_{\geq0}(R\dcate{Mod})$
	% segment-id: o014.li2.chapter8.bar-construction-b.q010
	\[ \epsilon_M: \mathrm{C}(\mathrm{Bar}(M)) \to M; \]
	this is a quasi-isomorphism because
	Corollary~\ref{prop:bar-exactness-G} ensures that $\epsilon_M$ is a
	quasi-isomorphism at the level of $\Bbbk\dcate{Mod}$. Explicitly,
	% segment-id: o014.li2.chapter8.bar-construction-b.q011
	\[ \mathrm{C}(\mathrm{Bar}(M))_n = L^{n+1}(M) = R^{\otimes (n+1)} \otimes M. \]
	According to Example~\ref{eg:monad-simplicial} and the description of
	$L\to\identity$ in Example~\ref{eg:comonad-change-ring} (namely,
	``multiplying in''), the map $L^{n+1}(M)\to L^n(M)$ is
	% segment-id: o014.li2.chapter8.bar-construction-b.q012
	\begin{multline*}
		\partial_n: r_0 \otimes \cdots \otimes r_n \otimes x \mapsto \\
		\sum_{k=0}^{n-1} (-1)^k \cdots \otimes r_k r_{k+1} \otimes \cdots \otimes x + (-1)^n r_0 \otimes \cdots \otimes r_{n-1} \otimes r_n x,
	\end{multline*}
	where $x\in M$ and $r_i\in R$. This formula still makes sense when
	$n=0$ ($r_0\otimes x\mapsto r_0x$), and thus gives the augmented version
	$\epsilon:\mathrm{C}(\mathrm{Bar}(M))\to M$, namely
	$\mathrm{C}^{\mathrm{aug}}(\mathrm{Bar}(M))$.

	Now consider the special case $M=R$. In this case
	% segment-id: o014.li2.chapter8.bar-construction-b.q013
	\begin{align*}
		\mathrm{C}(\mathrm{Bar}(R))_n & = R^{\otimes (n+2)} = R \otimes R^{\otimes n} \otimes R, \\
		\partial_n(r_0 \otimes \cdots \otimes r_{n+1}) & = \sum_{k=0}^n (-1)^k \cdots \otimes r_k r_{k+1} \otimes \cdots.
	\end{align*}

	Although $\mathrm{C}(\mathrm{Bar}(R))$ and its augmentation
	$\mathrm{C}^{\mathrm{aug}}(\mathrm{Bar}(R))$ are, by definition, chain
	complexes of left $R$-modules, the description above upgrades both to
	chain complexes of $(R,R)$-bimodules: $R$ acts by left multiplication on
	the first factor and by right multiplication on the last factor of
	$R\otimes R^{\otimes n}\otimes R$. This upgrade of structure is no
	coincidence. Since $\mathrm{C}^{\mathrm{aug}}(\mathrm{Bar}(\mathord\cdot))$
	is a functor, every endomorphism of $R$ as a left $R$-module (for example,
	right multiplication) lifts naturally to the entire bar construction.

	Comparison with Definition~\ref{def:bar-Hochschild} shows immediately
	that these are precisely the chain complexes used to define Hochschild
	homology and cohomology:
	% segment-id: o014.li2.chapter8.bar-construction-b.q014
	\[ \mathrm{C}(\mathrm{Bar}(R)) = \mathsf{B} R, \quad \mathrm{C}^{\mathrm{aug}}(\mathrm{Bar}(R)) = \mathsf{B}' R. \]
	Notice also that Lemma~\ref{prop:HH-bar-exactness} is just a special case
	of Theorem~\ref{prop:bar-null-homotopic}.

	Write $R^e:=R\otimes R^{\opp}$ and, according to
	Convention~\ref{con:Re-R-R}, identify left $R^e$-modules,
	$(R,R)$-bimodules, and right $R^e$-modules. From this viewpoint, the
	Hochschild homology $\HHm_\bullet(M)$ (respectively, cohomology
	$\HHm^\bullet(M)$) of a bimodule $M$ is a derived version of
	$M\dotimes{R^e}R$ (respectively, $\Hom_{R^e}(R,M)$): the bimodule $R$ is
	replaced by its bar resolution $\mathrm{C}(\mathrm{Bar}(R))$, after which
	homology (respectively, cohomology) is computed.

	This idea of replacement pervades the study of derived functors and
	derived categories, but the situation here is slightly different. In the
	framework of \CHref{sec:cplx} and \CHref{sec:triangulated-derived-cat},
	$R$ should be replaced by its flat resolution (respectively, projective
	resolution) as a bimodule. This yields $\Tor^{R^e}_\bullet(M,R)$
	(respectively, $\Ext^\bullet_{R^e}(R,M)$), or the corresponding object in
	the derived category. Example~\ref{eg:HH-Tor} (respectively,
	Example~\ref{eg:HH-Ext}) shows that if $R$ is flat (respectively,
	projective) as a $\Bbbk$-module, these objects are naturally isomorphic to
	$\HHm_\bullet(M)$ (respectively, $\HHm^\bullet(M)$); this need not hold
	in general.
\end{example}

% segment-id: o014.li2.chapter8.bar-construction-b.eg002
\begin{example}[The bar resolution and group homology]\label{eg:group-H-comonad}
	Let $\Gamma$ be a group. Continue the discussion of
	Example~\ref{eg:HH-comonad} for the group algebra $R:=\Bbbk[\Gamma]$; in
	this case, $M$ is the same as a left $\Gamma$-module in the sense of
	\S\ref{sec:G-mod}. Recall that
	$(\mathord\cdot)_{\Gamma}:\Gamma\dcate{Mod}\to\Bbbk\dcate{Mod}$ is the
	coinvariants functor. We claim that
	$\epsilon_M:\mathrm{C}(\mathrm{Bar}(M))\to M$ is a
	$(\mathord\cdot)_{\Gamma}$-acyclic resolution in $\Gamma\dcate{Mod}$
	(Convention~\ref{con:F-acyclic}). Its quasi-isomorphism property is
	already known. To see that every
	$\mathrm{Bar}(M)_n=\Bbbk[\Gamma]^{\otimes(n+1)}\otimes M$ is
	$(\mathord\cdot)_{\Gamma}$-acyclic, first use Shapiro's lemma
	(Theorem~\ref{prop:Shapiro}), which gives, for every $\Bbbk$-module $P$,
	% segment-id: o014.li2.chapter8.bar-construction-b.q015
	\[ \Hm_k(\Gamma, \Bbbk[\Gamma] \otimes P ) = \Hm_k\left( \Gamma, \iInd^{\Gamma}_{\{1\}}(P) \right) \simeq \Hm_k(\{1\}, P), \]
	and the right-hand side is always $0$ when $k>0$. Apply this with
	$P=\Bbbk[\Gamma]^{\otimes n}\otimes M$ to obtain the claim.

	Consequently, Corollary~\ref{prop:acyclic-resolution} implies
	$\Hm_n(\Gamma,M)\simeq
	\Hm_n\bigl(\mathrm{C}((\mathrm{Bar}(M)_{\Gamma}))\bigr)$; here the
	functor $(\mathord\cdot)_{\Gamma}$ may be applied either to the
	simplicial object or, equivalently, to the corresponding chain complex.
	This interprets group homology as comonad homology in the sense of the
	exercises in this chapter. More precisely, taking $\Gamma$-coinvariants
	of $\mathrm{Bar}(M)_n$ amounts to removing the first tensor factor in
	$\Bbbk[\Gamma]^{\otimes(n+1)}\otimes M$ via
	$\Bbbk[\Gamma]\twoheadrightarrow\Bbbk$, giving
	$\Bbbk[\Gamma]^{\otimes n}\otimes M$. Under this identification,
	$\mathrm{Bar}(M)_{n,\Gamma}\to\mathrm{Bar}(M)_{n-1,\Gamma}$ is given by
	% segment-id: o014.li2.chapter8.bar-construction-b.q016
	\begin{multline*}
		g_1 \otimes \cdots \otimes g_n \otimes x \mapsto g_2 \otimes \cdots \otimes g_n \otimes x + \sum_{k=1}^{n-1} (-1)^k \cdots \otimes g_k g_{k+1} \otimes \cdots \otimes x \\
		+ (-1)^n g_1 \otimes \cdots \otimes g_{n-1} \otimes g_n x,
	\end{multline*}
	where $g_i\in\Gamma$ and $x\in M$. Comparison with
	Proposition~\ref{prop:grouph-chain} shows immediately that
	$\mathrm{C}(\mathrm{Bar}(M))_{\Gamma}$ is precisely the standard chain
	complex computing $\Hm_n(\Gamma,M)$.
\end{example}
